\documentclass{article}
\usepackage{amsmath,amssymb,amsthm}
\usepackage{algorithm}
\usepackage{algpseudocode}
\usepackage{enumitem}
\usepackage{hyperref}
\usepackage{bm}
\newtheorem{theorem}{Theorem}
\newtheorem{lemma}{Lemma}
\newtheorem{assumption}{Assumption}
\newcommand{\E}{\mathbb{E}}
\newcommand{\Var}{\operatorname{Var}}
\newcommand{\R}{\mathbb{R}}
\newcommand{\norm}[1]{\left\|#1\right\|}
\newcommand{\ip}[2]{\left\langle #1,#2\right\rangle}
\newcommand{\grad}{\nabla}
\begin{document}

\section*{Algorithm Name}
\textbf{Stochastic Gradient Langevin Dynamics with Decaying Noise (SGLD‑D)}  

The method performs a stochastic gradient step and injects isotropic Gaussian noise whose variance decays with the iteration index.

\section*{Mathematical Formulation}
Let $f:\R^{d}\rightarrow\R$ be the objective.  
At iteration $k\ge 0$ the algorithm produces a sequence $\{x_k\}$ by  

\begin{equation}
\boxed{
x_{k+1}=x_{k}-\alpha_k\, g_k+\sqrt{2\alpha_k\,\sigma_k^{2}}\;\xi_k},
\qquad 
g_k:=\grad f(x_k)+\varepsilon_k,
\end{equation}
where  

\begin{itemize}[noitemsep,topsep=0pt]
\item $\alpha_k>0$ is the stepsize (learning rate);
\item $g_k$ is a (possibly stochastic) gradient estimator; $\varepsilon_k$ denotes its stochastic error, assumed unbiased;
\item $\xi_k\sim\mathcal N(0,I_d)$ is a standard Gaussian vector independent of the past;
\item $\sigma_k^{2}>0$ is a prescribed variance schedule, decreasing to $0$ (e.g. $\sigma_k^{2}=c/(k+1)^{\beta}$, $\beta>0$).
\end{itemize}

When $\sigma_k^{2}=0$ the method reduces to ordinary stochastic gradient descent (SGD).  

\section*{Pseudocode}
\begin{algorithm}[H]
\caption{SGLD‑D}
\begin{algorithmic}[1]
\State \textbf{Input:} initial point $x_0\in\R^{d}$, stepsize schedule $\{\alpha_k\}$, noise schedule $\{\sigma_k^{2}\}$, max iterations $K$.
\For{$k=0,\dots,K-1$}
    \State Sample a stochastic gradient $g_k$ such that $\E[g_k\mid x_k]=\grad f(x_k)$.
    \State Sample $\xi_k\sim\mathcal N(0,I_d)$.
    \State $x_{k+1}\gets x_k-\alpha_k g_k+\sqrt{2\alpha_k\sigma_k^{2}}\;\xi_k$.
\EndFor
\State \textbf{Return} $x_K$ (or the averaged iterate $\bar x_K=\frac1K\sum_{k=0}^{K-1}x_k$).
\end{algorithmic}
\end{algorithm}

\section*{Dry Run Example}
Consider the one‑dimensional quadratic $f(w)=(w-3)^2$.  
Gradient: $\grad f(w)=2(w-3)$.  
Set deterministic gradients (no stochastic error), stepsize $\alpha=0.1$, and a simple variance schedule $\sigma_k^{2}=0$ (so the algorithm coincides with gradient descent).  

\begin{center}
\begin{tabular}{c|c|c|c}
$k$ & $w_k$ & $g_k=2(w_k-3)$ & $w_{k+1}=w_k-\alpha g_k$ \\ \hline
0 & $0$ & $-6$ & $0-0.1(-6)=0.6$ \\
1 & $0.6$ & $-4.8$ & $0.6-0.1(-4.8)=1.08$ \\
2 & $1.08$ & $-3.84$ & $1.08-0.1(-3.84)=1.464$ \\
\end{tabular}
\end{center}
Corresponding objective values: $f(0)=9$, $f(0.6)=5.76$, $f(1.08)=3.6864$, $f(1.464)=2.3529$.  

The example shows the descent direction and the reduction of $f$ over three iterations.

\newpage
\section*{Assumptions}
\begin{assumption}[Smoothness]
$f$ is $L$‑smooth, i.e. for all $x,y\in\R^{d}$,
\[
f(y)\le f(x)+\ip{\grad f(x)}{y-x}+\frac{L}{2}\norm{y-x}^{2}.
\tag{A1}
\]
\end{assumption}
\begin{assumption}[Bounded Gradient Variance]
The stochastic gradient error satisfies
\[
\E[\varepsilon_k\mid x_k]=0,\qquad 
\E[\norm{\varepsilon_k}^{2}\mid x_k]\le \sigma_{g}^{2},
\tag{A2}
\]
for a constant $\sigma_{g}^{2}\ge0$.
\end{assumption}
\begin{assumption}[Noise Schedule]
The injected noise variance is non‑increasing and satisfies
\[
\sigma_k^{2}\le \frac{c}{(k+1)^{\beta}},\qquad c>0,\;\beta>0.
\tag{A3}
\]
\end{assumption}
\begin{assumption}[Stepsize]
The stepsizes are constant or slowly decreasing such that
\[
0<\alpha_k\le \alpha_{\max}<\frac{1}{L},\qquad 
\sum_{k=0}^{\infty}\alpha_k=\infty,\;\;
\sum_{k=0}^{\infty}\alpha_k^{2}<\infty .
\tag{A4}
\]
\end{assumption}

\section*{Main Convergence Theorem}
\begin{theorem}[Convergence of SGLD‑D]\label{thm:main}
Assume (A1)–(A4) hold. Let $x^{*}$ be a (global) minimizer of $f$ and define $F_{*}=f(x^{*})$. For any $K\ge1$,
\[
\frac{1}{K}\sum_{k=0}^{K-1}\E\!\big[\norm{\grad f(x_k)}^{2}\big]
\;\le\;
\frac{2\big(f(x_0)-F_{*}\big)}{\alpha_{\min}K}
\;+\;
L\big(\sigma_{g}^{2}+2\sigma_{\max}^{2}\big),
\tag{1}
\]
where $\alpha_{\min}:=\min_{k< K}\alpha_k$ and $\sigma_{\max}^{2}:=\max_{k< K}\sigma_k^{2}$.  
Consequently, if $\sigma_k^{2}=O(k^{-\beta})$ with $\beta>0$, the right‑hand side decays as $O(1/K)$, establishing that the average squared gradient norm tends to zero.
\end{theorem}

\section*{Theoretical Properties}
\begin{enumerate}[noitemsep,topsep=0pt]
\item \textbf{Stability.} Because $\alpha_k<L^{-1}$, the deterministic part of the update is a contraction in expectation (see Lemma~\ref{lem:descent}).
\item \textbf{Hyperparameter Sensitivity.} The bound (1) shows a linear dependence on $L$, the variance constants $\sigma_{g}^{2}$ and $\sigma_{\max}^{2}$, and an inverse dependence on the stepsize lower bound $\alpha_{\min}$. Smaller stepsizes improve the bias term but enlarge the variance term through the factor $L$.
\item \textbf{Comparison to SGD.} Setting $\sigma_k^{2}=0$ recovers the classical SGD bound 
$\frac{1}{K}\sum_{k}\E\!\big[\norm{\grad f(x_k)}^{2}\big]\le
\frac{2(f(x_0)-F_{*})}{\alpha_{\min}K}+L\sigma_{g}^{2}$, i.e. the additional $2L\sigma_{\max}^{2}$ term is the price paid for exploration.
\item \textbf{Stationary Distribution.} When $\sigma_k^{2}\equiv\sigma^{2}>0$ the iterates asymptotically sample from a distribution whose density is proportional to $\exp\!\big(-f(x)/\sigma^{2}\big)$ (standard SGLD result). The decaying schedule forces the temperature to zero, thus concentrating the distribution around stationary points.
\end{enumerate}

\section*{Discussion}
The theorem guarantees that, despite the presence of diminishing stochastic noise, the algorithm does not stall at high‑variance points: the expected squared gradient norm diminishes at the classic $O(1/K)$ rate for non‑convex smooth objectives. The proof follows the classical descent‑lemma argument augmented with explicit handling of the injected Gaussian noise.

\newpage
\section*{Preliminaries (Definitions and Lemmas)}
\begin{lemma}[Descent Lemma under $L$‑smoothness]\label{lem:descent}
For any $x\in\R^{d}$ and any vector $u\in\R^{d}$,
\[
f(x-u)\le f(x)-\ip{\grad f(x)}{u}+\frac{L}{2}\norm{u}^{2}.
\tag{2}
\]
\end{lemma}
\begin{proof}
Directly from (A1) with $y=x-u$.
\end{proof}

\section*{Main Proof (Step‑by‑Step Derivation)}
We prove Theorem~\ref{thm:main}.  
All expectations are taken with respect to the randomness up to iteration $k$ unless otherwise noted.

\noindent\textbf{[Step 1]} Apply Lemma~\ref{lem:descent} with $u=\alpha_k g_k-\sqrt{2\alpha_k\sigma_k^{2}}\xi_k$:
\[
f(x_{k+1})\le f(x_k)-\ip{\grad f(x_k)}{\alpha_k g_k}
+\frac{L}{2}\norm{\alpha_k g_k-\sqrt{2\alpha_k\sigma_k^{2}}\xi_k}^{2}.
\tag{3}
\]

\noindent\textbf{[Step 2]} Expand the squared norm in (3):
\[
\begin{aligned}
\norm{\alpha_k g_k-\sqrt{2\alpha_k\sigma_k^{2}}\xi_k}^{2}
=&\alpha_k^{2}\norm{g_k}^{2}
+2\alpha_k\sigma_k^{2}\norm{\xi_k}^{2}
-2\alpha_k^{3/2}\sqrt{2\sigma_k^{2}}\;\ip{g_k}{\xi_k}.
\end{aligned}
\tag{4}
\]

\noindent\textbf{[Step 3]} Take conditional expectation $\E[\cdot\mid x_k]$.
Because $\xi_k$ is independent of $g_k$ and $\E[\xi_k]=0$, the cross term vanishes:
\[
\E\!\big[\ip{g_k}{\xi_k}\mid x_k\big]=0.
\]
Moreover $\E[\norm{\xi_k}^{2}]=d$. Hence
\[
\E\!\big[\norm{\alpha_k g_k-\sqrt{2\alpha_k\sigma_k^{2}}\xi_k}^{2}\mid x_k\big]
=
\alpha_k^{2}\E\!\big[\norm{g_k}^{2}\mid x_k\big]
+2\alpha_k\sigma_k^{2}d .
\tag{5}
\]

\noindent\textbf{[Step 4]} Decompose $g_k=\grad f(x_k)+\varepsilon_k$ and use unbiasedness:
\[
\E\!\big[\norm{g_k}^{2}\mid x_k\big]
= \norm{\grad f(x_k)}^{2}
+\E\!\big[\norm{\varepsilon_k}^{2}\mid x_k\big]
\le \norm{\grad f(x_k)}^{2}+\sigma_{g}^{2},
\tag{6}
\]
where (A2) was used.

\noindent\textbf{[Step 5]} Combine (3)–(6) and take total expectation:
\[
\begin{aligned}
\E[f(x_{k+1})]
\le &\E\!\big[f(x_k)-\alpha_k\ip{\grad f(x_k)}{g_k}\big]
+\frac{L}{2}\E\!\big[\alpha_k^{2}\norm{g_k}^{2}+2\alpha_k\sigma_k^{2}d\big] \\
\le &\E\!\big[f(x_k)-\alpha_k\norm{\grad f(x_k)}^{2}\big]
+\frac{L\alpha_k^{2}}{2}\E\!\big[\norm{\grad f(x_k)}^{2}+\sigma_{g}^{2}\big]
+\;L\alpha_k\sigma_k^{2}d .
\end{aligned}
\tag{7}
\]

\noindent\textbf{[Step 6]} Rearrange (7) to isolate $\norm{\grad f(x_k)}^{2}$:
\[
\big(\alpha_k-\tfrac{L}{2}\alpha_k^{2}\big)\E\!\big[\norm{\grad f(x_k)}^{2}\big]
\le \E[f(x_k)-f(x_{k+1})]
+\tfrac{L}{2}\alpha_k^{2}\sigma_{g}^{2}+L\alpha_k\sigma_k^{2}d .
\tag{8}
\]

\noindent\textbf{[Step 7]} Since $\alpha_k\le\alpha_{\max}<1/L$, we have $\alpha_k-\frac{L}{2}\alpha_k^{2}\ge\frac{\alpha_k}{2}$. Thus
\[
\frac{\alpha_k}{2}\E\!\big[\norm{\grad f(x_k)}^{2}\big]
\le \E[f(x_k)-f(x_{k+1})]
+\tfrac{L}{2}\alpha_k^{2}\sigma_{g}^{2}+L\alpha_k\sigma_k^{2}d .
\tag{9}
\]

\noindent\textbf{[Step 8]} Sum (9) over $k=0,\dots,K-1$:
\[
\frac{1}{2}\sum_{k=0}^{K-1}\alpha_k\E\!\big[\norm{\grad f(x_k)}^{2}\big]
\le f(x_0)-\E[f(x_K)]
+\tfrac{L}{2}\sigma_{g}^{2}\sum_{k=0}^{K-1}\alpha_k^{2}
+L d\sum_{k=0}^{K-1}\alpha_k\sigma_k^{2}.
\tag{10}
\]

\noindent\textbf{[Step 9]} Drop the non‑negative term $-\E[f(x_K)]\le 0$ and use $f(x_K)\ge F_{*}$:
\[
\frac{1}{2}\sum_{k=0}^{K-1}\alpha_k\E\!\big[\norm{\grad f(x_k)}^{2}\big]
\le f(x_0)-F_{*}
+\tfrac{L}{2}\sigma_{g}^{2}\sum_{k=0}^{K-1}\alpha_k^{2}
+L d\sum_{k=0}^{K-1}\alpha_k\sigma_k^{2}.
\tag{11}
\]

\noindent\textbf{[Step 10]} Divide (11) by $\sum_{k=0}^{K-1}\alpha_k\ge K\alpha_{\min}$:
\[
\frac{1}{K}\sum_{k=0}^{K-1}\E\!\big[\norm{\grad f(x_k)}^{2}\big]
\le
\frac{2\big(f(x_0)-F_{*}\big)}{\alpha_{\min}K}
+L\sigma_{g}^{2}\frac{\sum_{k=0}^{K-1}\alpha_k^{2}}{K\alpha_{\min}}
+2L d\frac{\sum_{k=0}^{K-1}\alpha_k\sigma_k^{2}}{K\alpha_{\min}}.
\tag{12}
\]

\noindent\textbf{[Step 11]} Because $\alpha_k\le\alpha_{\max}$ and $\sigma_k^{2}\le\sigma_{\max}^{2}$,
\[
\frac{\sum_{k=0}^{K-1}\alpha_k^{2}}{K\alpha_{\min}}\le\alpha_{\max},
\qquad
\frac{\sum_{k=0}^{K-1}\alpha_k\sigma_k^{2}}{K\alpha_{\min}}\le\sigma_{\max}^{2}.
\]
Insert these bounds into (12) and absorb the dimension constant $d$ into $\sigma_{\max}^{2}$ (define $\tilde\sigma_{\max}^{2}=d\sigma_{\max}^{2}$). We obtain
\[
\frac{1}{K}\sum_{k=0}^{K-1}\E\!\big[\norm{\grad f(x_k)}^{2}\big]
\le
\frac{2\big(f(x_0)-F_{*}\big)}{\alpha_{\min}K}
+L\big(\sigma_{g}^{2}+2\tilde\sigma_{\max}^{2}\big),
\]
which is exactly the bound (1). ∎

\section*{Rate Derivation (With Constants)}
If the variance schedule follows $\sigma_k^{2}=c/(k+1)^{\beta}$ with $\beta>0$, then  
\[
\sum_{k=0}^{K-1}\alpha_k\sigma_k^{2}\le \alpha_{\max}c\sum_{k=1}^{K}\frac{1}{k^{\beta}}
=
\begin{cases}
O(1) & \beta>1,\\[2mm]
O(\log K) & \beta=1,\\[2mm]
O(K^{1-\beta}) & 0<\beta<1.
\end{cases}
\]
Since this term is divided by $K$ in (12), its contribution is $O(1/K)$ for any $\beta>0$.  
Thus the dominant term is the $O(1/K)$ bias term, confirming the $O(1/K)$ convergence rate for the average squared gradient norm.

\section*{Special Cases}
\begin{itemize}[noitemsep,topsep=0pt]
\item \textbf{Pure SGD ($\sigma_k^{2}=0$).} The bound reduces to the classical non‑convex SGD result.
\item \textbf{Constant temperature ($\sigma_k^{2}\equiv\sigma^{2}>0$).} The second variance term becomes $L\sigma^{2}$, yielding a steady‑state bias that matches the known SGLD stationary distribution error.
\item \textbf{Deterministic gradients ($\varepsilon_k\equiv0$).} Only the injected noise contributes; the bound simplifies to $\frac{2(f(x_0)-F_{*})}{\alpha_{\min}K}+2L\sigma_{\max}^{2}$.
\end{itemize}

\end{document}